\scanpage{27}
% Pencil foliation: 26, lower left.
dont la formation est compatible avec chgt
de base (il suffit donc de le construire
sur $\operatorname{Spec}\mathbb{Z}$), avec \deleted{des} isom
\[
 T_\ell(\unsure{\mathbb{Q}|\mathbb{Z}}(\unsure{-1|1}))
 \simeq \unsure{\mathbb{Q}|\mathbb{Z}}_\ell(\unsure{-1|1})
 =\underbrace{T_\ell(\mathbb{G}_{\mathrm m})^{-1}}_{
 \substack{\text{Module de Tate}\\\text{\inserted{inverse} de }\mathbb{G}_{\mathrm m}}}
\]
et le cas échéant ($X$ sur $\mathbb{Q}$).

On peut définir $\unsure{\mathbb{Q}|\mathbb{Z}}(-1)$ comme
$R^2f_*(\deleted{\unsure{\mathbb{Z}}}\mathbf1_{\mathbb{P}^1_X})$,
où $f:\mathbb{P}^1_X\longrightarrow X$ est la
projection canonique. Posant
\[
 \unsure{\mathbb{Q}|\mathbb{Z}}(-n)=\unsure{\mathbb{Q}|\mathbb{Z}}(-1)^{\otimes n}
 \qquad\text{pour }n\ge0
\]
on peut prouver, à l'aide des
axiomes déjà posés, que si $f:X\unsure{-|\longrightarrow}S$
est \inserted{lisse} projectif à fibres \inserted{géom.} connexes non vides,
\deleted{partout} de dim relative $d$, alors
\[
 R^{2d}f_*(\mathbf1_X)\simeq\mathbb{Z}_S(-d)
\]
et si on enlève l'hyp. ``projectif'' mais
suppose $f$ quasi-projectif, on trouve encore
\[
 R^{2d}f_!(\mathbf1_X)\simeq\mathbb{Z}_S(-d)
\]
